\section{Problem formulation}
\label{sec:problem}

\subsection{Forecasts, hits and the quantities to repair}
\label{sec:problem-setting}

A portfolio has a simple return $y_t$ on trading day $t$, and its loss is $L_t = -y_t$, a positive number on a day of loss. Let $\mathcal{F}_{t-1}$ be the information available before day $t$ opens. A \emph{model} is a map that, for every tail level $a \in (0, \tfrac12]$, returns a one-day value at risk $V_t(a)$ and an expected shortfall $E_t(a)$, both positive losses measured in decimals of the portfolio value, computed from $\mathcal{F}_{t-1}$ only. The pair satisfies $E_t(a) \ge V_t(a)$ and, for every day, $V_t(a)$ is non-increasing in $a$: a smaller tail level asks for a larger loss. For an ideal model, $V_t(a)$ is the $(1-a)$-quantile of the conditional law of $L_t$ and $E_t(a)$ is the mean loss beyond it,
\begin{equation}
V_t(a) = \inf\{x : \Prob(L_t \le x \mid \mathcal{F}_{t-1}) \ge 1-a\}, \qquad E_t(a) = \E\bigl[L_t \mid L_t \ge V_t(a), \mathcal{F}_{t-1}\bigr].
\label{eq:var-es}
\end{equation}
The forecaster is asked for the pair at one target level $\alpha$ (here 0.01 and 0.025). Day $t$ is a \emph{hit} when the loss exceeds the forecast,
\begin{equation}
e_t(a) = \indic\{L_t > V_t(a)\},
\label{eq:hit}
\end{equation}
and the forecasts are \emph{calibrated} when the hits satisfy two properties in the sense of \citet{christoffersen1998}: unconditional coverage, $\Prob(e_t = 1) = \alpha$, and independence, $e_t$ unpredictable from the past hits. Together they are the statement that $e_t - \alpha$ is a martingale difference with respect to the information that includes the past hits. The ES is judged separately, through a score for which the pair $(V_t(\alpha), E_t(\alpha))$ of \cref{eq:var-es} is the unique minimiser of the expected loss. In the return convention of the engine (returns $y_t = -L_t$, forecasts $v_t = -V_t < 0$ and $\epsilon_t = -E_t < 0$) the score of \citet{pattonziegelchen2019}, a member of the family of \citet{fissler2016}, is
\begin{equation}
\mathrm{FZ0}(y_t, v_t, \epsilon_t; \alpha) = -\frac{1}{\alpha\,\epsilon_t}\,\indic\{y_t \le v_t\}\,(v_t - y_t) + \frac{v_t}{\epsilon_t} + \log(-\epsilon_t) - 1,
\label{eq:fz0}
\end{equation}
which is a score: only differences between models are meaningful, and the score of one model can be negative.

\subsection{The recalibration problem}
\label{sec:problem-recal}

Given a model whose estimator, window and refit schedule are fixed (a constraint of a bank's production environment, where a replacement of the model passes through a model-risk process), find a rule that chooses, on each day $t$, a working level $a_t \in [\ell, u]$ from the past hits only, so that the pair $(V_t(a_t), E_t(a_t))$ is delivered in place of $(V_t(\alpha), E_t(\alpha))$ and the three properties improve: the long-run frequency of hits approaches $\alpha$ under a change of regime, the independence of the hits is not damaged, and the joint score of \cref{eq:fz0} does not worsen by more than a stated tolerance. The tolerance used here is a relative change of the mean score of 5\%, the same figure for every model and window; it is a decision of the product and not a statistical bound. The ES is not recalibrated by a rule of its own: it is whatever the model gives at the working level, and the question of how well it is calibrated after the change is measured by the second Acerbi--Szekely statistic \citep{acerbiszekely2014},
\begin{equation}
Z_2 = 1 - \frac{1}{T\alpha}\sum_{t=1}^{T} \frac{L_t\, e_t}{E_t},
\label{eq:z2}
\end{equation}
which has expectation zero under a correct pair and is negative when the ES is understated.

\subsection{Assumptions}
\label{sec:problem-assumptions}

The analysis uses five assumptions, each stated where it enters.
\begin{enumerate}
\item The forecasts $V_t(a)$ and $E_t(a)$ use $\mathcal{F}_{t-1}$ only. The wrapper cannot check this; the experiments check it by construction, since every model in the ladder slices its window before day $t$.
\item For every day, $V_t(\cdot)$ is continuous and strictly decreasing on $[\ell, u]$ with $V_t(a) \to \infty$ as $a \to 0$, so that every level in the band gives a finite VaR. This fails for historical simulation below the reciprocal of its window, where the VaR is flat; the band's lower edge of 0.001 lies in that flat region for a window of 250 days, which the experiments report through the count of days at the band edge.
\item The loss $L_t$ may follow any process: the identity of \cref{sec:theory} makes no assumption on it. The fixed-point analysis assumes, further, a stationary scale family (\cref{prop:fixed-point}).
\item The model is evaluated at each day in time order, and a day for which the model gives no usable forecast keeps the forecast of the previous day, as the ladder does; the number of such days is reported (\tenFallback{} days in all, over the \tenIndRunsAll{} wrapped runs).
\item The hits are observed at the end of each day without delay. A delay of $k$ days in the confirmation of a hit (the case of a bank that books trades late) is out of scope.
\end{enumerate}
The notation is collected in \cref{tab:notation}; the symbols of the tests and of the experiments are defined where they appear.

\begin{table}[htbp]
\centering
\small
\caption{Notation of the manuscript: each symbol with its meaning and its unit; the symbols of the back-tests and scores are those defined where they are first used in the experimental design.}
\label{tab:notation}
\begin{adjustbox}{max width=\textwidth}
\begin{tabular}{lll}
\toprule
Symbol & Meaning & Unit \\
\midrule
$y_t$, $L_t$ & \relax Return and loss of day $t$ ($L_t = -y_t$) & decimal \\
$\alpha$ & Target tail level of the VaR & probability \\
$a_t$ & \relax Working tail level on day $t$ & probability \\
$\gamma$ & Step of the level update & probability \\
$[\ell, u]$ & Clipping band of the working level & probability \\
$V_t(a)$, $E_t(a)$ & \relax VaR and ES of the model at level $a$ (positive loss) & decimal \\
$e_t$ & \relax Hit indicator $\indic\{L_t > V_t(a_t)\}$ & 0 or 1 \\
$T$ & Number of forecast days & days \\
$C_T$ & Total amount the band removes from the level & probability \\
$c$ & Ratio of the true scale to the scale the model assumes & 1 \\
$F$, $f$, $q_\nu$ & \relax Distribution, density and quantile function of the unit-variance Student-$t$ & 1 \\
$a^*$ & Fixed point of the level update & probability \\
$S_t$, $\theta_t$ & \relax Score $L_t/V_t$ and the scale applied to the VaR & 1 \\
FZ0 & Mean Fissler--Ziegel score of the VaR and ES pair & score \\
$Z_2$ & Acerbi--Szekely ES statistic & 1 \\
\bottomrule
\end{tabular}
\end{adjustbox}
\end{table}

